\subsection{Infinitely changed exponents which are units almost everywhere}

\begin{statement}{A.65}{All almost-everywhere-unit Selmer exponents are integers}{Proved}
Let $E/\Q$ be any elliptic curve, let $N\ge2$, put
$U=E-E[N]$, and let $R\in E(\Q)$ have infinite order. Define
\[
 \mathcal U_{\rm fin}=
 \{a=(a_p)_p\in\widehat\Z:
           a_p\in\Z_p^\times\text{ for all but finitely many }p\}.
 \tag{A65.1}
\]
Among these exponents, the proper section with compact Kummer
class $a\kappa(R)$ has a geometrically abelian Selmer lift if and
only if $a\in\Z\setminus\{0\}$. At every such integer $a=n$
the abelian Selmer fibre is the single point section of $nR$.
The same exact set of exponents occurs for Selmer lifts to the
intrinsic metabelian quotient, without an assertion about the
metabelian fibres.

This assertion requires neither complex multiplication, rank one,
nor finite Tate--Shafarevich group. It includes exponents changed
at infinitely many coefficient primes. Only local point-or-cusp
geometry at every finite place is assumed.

We also prove the following erasure strengthening over any number
field $k$: for every $u\in\widehat\Z^\times$ and every nonempty
finite set $\mathcal P$ of coefficient primes, the compact class
\[
 u(1-e_{\mathcal P})\kappa(R),\qquad R\in E(k)
                 \text{ of infinite order},
 \tag{A65.2}
\]
has no geometrically abelian Selmer lift on $E-E[N]$.
\end{statement}
\begin{proof}
We first treat $a=u\in\widehat\Z^\times$ over $\Q$, allowing any
nontorsion rational point $R$. The full local Tate Kummer map
is injective on $E(\Q_q)$ for every finite prime $q$. This
local group is compact and profinite, and its continuous
$\widehat\Z$-module structure therefore defines
\[
 Q_q=[u]R\in E(\Q_q).
 \tag{A65.3}
\]
Multiplication by $u$ is an automorphism. Since $R$ remains
nontorsion over $\Q_q$, so does $Q_q$. In particular
$Q_q\notin E[N]$, including at the bad primes. Thus a
hypothetical abelian Selmer lift has the interior point $Q_q$
as its forced local centre at \emph{every} finite place.

There is a nonconstant rational unit
$f\in\Q[U]^\times$. Indeed, over a finite Galois extension
$L/\Q$ splitting $E[N]$, take a nonzero $d\in E[N](L)$ and
a function with divisor $N(d)-N(O)$, and take its field norm.
The resulting divisor has a nonzero negative coefficient at
$O$ and an effective nonzero part supported away from $O$;
it cannot be zero. This proves both rationality and
nonconstancy of the norm.

Choose a finite set $S$ containing the primes of bad reduction,
the primes dividing $N$, and the vertical divisor components
of $f$ on a smooth model away from those primes. Enlarge $S$
so that $[N]R$ does not reduce to $O$ outside $S$. Such an
enlargement is finite: $[N]R\ne O$ is a fixed rational point
and the two sections $[N]R$ and $O$ of the proper model
intersect at only finitely many primes. For $q\notin S$,
multiplication by $u$ is an automorphism also on the finite
reduction group. Hence
\[
 [N]\overline{Q_q}=0
       \quad\Longleftrightarrow\quad [N]\overline R=0,
\]
and neither condition holds. The point $Q_q$ therefore
reduces away from the divisor of $f$, giving
\[
 f(Q_q)\in\Z_q^\times\qquad(q\notin S).
 \tag{A65.4}
\]

The homomorphism induced by $f:U\to\mathbf G_m$ factors
through the geometric abelian quotient. Applying it to
the proposed global section gives a single class
\[
 \beta\in H^1(\Q,\widehat\Z(1))
       =\varprojlim_m\Q^\times/(\Q^\times)^m.
 \tag{A65.5}
\]
Its localization is the ordinary Kummer class of $f(Q_q)$.
Consequently its valuation at every $q$ is an ordinary
integer $n_q=v_q(f(Q_q))$, diagonally embedded in
$\widehat\Z$, and $n_q=0$ outside $S$.
These statements force $\beta$ to be the Kummer class of
an actual rational number. To see the completion issue
explicitly, put $r_0=\prod_{q\in S}q^{n_q}$. The class
$\beta/\kappa(r_0)$ has zero valuation at every prime.
At any finite level $m$, its representative in
$\Q^\times/(\Q^\times)^m$ consequently has all prime
valuations divisible by $m$, and is represented by a
sign. The inverse limit of these sign images is the
actual group $\{\pm1\}$. Thus
\[
 \beta=\kappa(r),\qquad r=\pm r_0\in\Q^\times.
 \tag{A65.6}
\]
This uses the entire completed class, not an assumption
that a general element of the completion has finite
support.

The full local multiplicative Kummer map is injective
on $\Q_q^\times$: in the decomposition into valuation,
finite roots of unity, and principal units, no nontrivial
element is divisible by every positive integer.
Equations (A65.5)--(A65.6) therefore give
\[
 f(Q_q)=r\quad\text{for all finite }q.
 \tag{A65.7}
\]
In particular all compact local points lie in the fixed
finite reduced $\Q$-scheme $Z=(f^{-1}(r))_{\rm red}\subset E$.

We spell out why this finite support identifies the
\emph{global proper section}. The geometric fundamental
group of $E$ is the full Tate module $T E$, so
$u\kappa(R)$ is an actual section of $\pi_1(E)$.
Take the cofiltered system of its geometrically connected
finite \'{e}tale section-neighbourhoods $E_i\to E$.
Every $E_i$ is a smooth proper genus-one curve, since
the covering is \'{e}tale. At each finite place, the
fixed-coset description of a lifted point section gives
a point of $E_i(\Q_q)$ above $Q_q$, hence in the finite
scheme
\[
 Z_i=(E_i\times_E Z)_{\rm red}.
\]
The real restriction of the compact Kummer section is
geometric as well: $u$ is odd, so its Kummer class
is that of the component of $R$ in $\pi_0(E(\R))$,
a group of exponent at most two.
The same fixed-coset argument lifts a
real point to each section-neighbourhood. Thus the
lifted section supplies a modified adelic point of
$E_i$ surviving finite descent by
\citet[Theorem~2.1]{AMHarariStix2012}. All its finite
coordinates belong to $Z_i$.
\citet[Theorem~8.2]{AMStoll2007}, with the finite-place
interpretation of the erratum \citep{AMStollErrata2017}, now gives
$Z_i(\Q)\ne\varnothing$.

These are finite nonempty sets, and a common refinement
supplies compatible points for any finite collection
of neighbourhoods. Compactness gives a point in their
inverse limit. Equivalently, the section fixes a
compatible pro-point above some $z\in Z(\Q)$; its
stabilizer is the point decomposition group at $z$.
There is no boundary inertia in this proper argument.
The proper section is therefore the rational point
section of $z$, by precisely the finite-support
neighbourhood argument of A.2, here with every
neighbourhood already of genus one. Merely identifying
the original adelic tuple would not replace this step.

Global Kummer injectivity
$E(\Q)^\wedge\hookrightarrow H^1(\Q,T E)$ now identifies
$uR$ with $z$ in the completion of the finitely generated
Mordell--Weil group. This implies $u\in\Z$ even when
$R$ is not primitive. For clarity, write the free
coordinates of $R$ and $z$ as
$(r_1,\ldots,r_t)$ and $(z_1,\ldots,z_t)$, choose
integers $b_i$ with $\sum_i b_i r_i=d>0$, and obtain
\[
 du=\sum_i b_i z_i\in\Z.
\]
Thus $u$ is one rational number lying in every $\Z_p$,
and is an integer. An integer in every
$\Z_p^\times$ is $1$ or $-1$. This proves the unit
case. Applying it to $mR$ proves the same conclusion
for $a=mu$, $m\in\Z_{>0}$: necessarily $a=\pm m$.

We next prove (A65.2). Pass to one finite extension
$L/k$ splitting $E[N]$. Choose an auxiliary prime
$\lambda\nmid N$ and the entire finite family
\[
 f_d\in L(E)^\times,\qquad
 \operatorname{div}(f_d)=N(d)-N(O),
 \qquad 0\ne d\in E[N](L).
 \tag{A65.8}
\]
Each function applied to the putative section gives
one global mod-$\lambda$ Kummer class. Choose a finite
set $S$ containing their representative supports,
bad reduction, vertical components, and the primes
above $N\lambda M$, where $M=\prod_{p\in\mathcal P}p$.
Outside $S$, all torsion sections are disjoint and
every $f_d$ has exactly the indicated horizontal divisor.

At a place $v\notin S$, write
\[
 R=A_v+T_v,\qquad
 A_v=(1-e_{\mathcal P})R,\quad
 T_v=e_{\mathcal P}R.
 \tag{A65.9}
\]
The $\mathcal P$-primary factor of $E(L_v)$ is finite,
since its open formal subgroup is pro-residue-characteristic
and that characteristic is not in $\mathcal P$.
Thus $T_v$ is finite torsion of order prime to the
residue characteristic. The forced local point is now
$Q_v=[u]A_v$. It is nontorsion, so lies in $U(L_v)$,
and the global unit characters give
\[
 v(f_d(Q_v))\equiv0\pmod\lambda
                 \quad(v\notin S,\ d\ne O).
 \tag{A65.10}
\]

Consider the distinct rational points $R_j=[M^j]R$,
$j\ge0$. If $R_j$ reduces into $E[N]$, the primary
decomposition of the finite reduction group gives
$N\overline{A_v}=0$ and
$N[M^j]\overline{T_v}=0$: multiplication by $M^j$
is invertible on the complementary primary part.
Reduction is injective on prime-to-residue-characteristic
torsion. Since $E[N]$ is split, there are therefore
$d_0,d_1\in E[N](L)$ such that
\[
 \overline{A_v}=\overline{d_0},\qquad
 [M^j]T_v=d_1,\qquad
 R_j-d'=[M^j](A_v-d_0),\quad
 d'=[M^j]d_0+d_1.
 \tag{A65.11}
\]
The positive formal distance of $A_v$ from $d_0$ is
preserved by multiplication by $M^j$. It is also
preserved by $u$: on the formal group this action
factors through the residue-characteristic component
of $u$, which is a unit, and both it and its inverse
preserve every formal filtration step. Moreover
$[u]d_0$ is the split torsion point
$[u\bmod N]d_0$. Hence
\[
 Q_v-[u]d_0=[u](A_v-d_0)
 \tag{A65.12}
\]
has exactly the same formal distance, say $h>0$.
Apply (A65.10) to $f_{[u]d_0}$ if $d_0\ne O$,
or to any $f_e$ if $d_0=O$. The valuation is
$Nh$ or $-Nh$, respectively, so $\lambda\mid h$.
Equation (A65.11) then implies, for every fixed
$d\ne O$ and every $j$,
\[
 v(f_d(R_j))\equiv0\pmod\lambda\qquad(v\notin S).
 \tag{A65.13}
\]
If $R_j$ does not reduce to a zero or pole, its
valuation is zero, which covers the remaining case.
This is why the full family in (A65.8), rather than
one arbitrarily chosen function, is used before
choosing the final cover.

The group of classes in
$L^\times/(L^\times)^\lambda$ whose valuations outside
$S$ are divisible by $\lambda$ is finite, by the
finiteness of the $S$-ideal-class group and finite
generation of the $S$-unit group. Fix $d\ne O$.
Equation (A65.13) makes all the $R_j$ lift to a
finite family of twists
$z^\lambda=c^{-1}f_d$. Each smooth projective
normalization is geometrically connected and has
genus $\lambda$: its degree-$\lambda$ map to $E$
ramifies fully at exactly $d$ and $O$, because
$\lambda\nmid N$. Faltings's theorem makes every
member's set of $L$-points finite
\citep[Satz~7, p.~365, and the correction]{AMFaltings1983,AMFaltings1984Erratum}.
This contradicts the infinitely many distinct
$R_j$, proving (A65.2). Notice that multiplying
$R$ by a nonzero integer is harmless, whereas
$R=O$ is explicitly excluded from this argument.

Finally, every $a\in\mathcal U_{\rm fin}$ factors
as follows. Its zero-component set $\mathcal P$
is finite. At the finitely many nonzero nonunit
components, put their valuations into
$m=\prod_{p\notin\mathcal P}p^{v_p(a_p)}>0$.
For $p\notin\mathcal P$ set $u_p=a_p/m$, and
choose any unit $u_p$ at the zero components.
Then
\[
 a=mu(1-e_{\mathcal P}),\qquad
 u\in\widehat\Z^\times.
 \tag{A65.14}
\]
If $\mathcal P$ is nonempty, (A65.2) applied to
$mR$ forbids a lift. Otherwise the rational-unit
argument gives $u=\pm1$, so $a=\pm m$.
Conversely every nonzero integer gives the actual
point section of $nR\in U(\Q)$, and A.64 makes
its abelian Selmer fibre unique. Any metabelian
Selmer lift projects to an abelian one, while
actual point sections lift to the full group.
This proves the asserted equality of exponent sets.
\end{proof}

The new family is not restricted to finite changes of an
integer. For example, take $a_2=1$, and for odd primes
take $a_p=1$ when $p\equiv1\pmod4$ and $a_p=-1$ when
$p\equiv3\pmod4$. This is a unit in every coefficient
component, but is neither of the diagonal integers
$1,-1$; A.65 excludes its Selmer lift. For the
least-conductor cubic, the still undecided nonintegral
exponents must have nonunit components at infinitely
many primes; A.63 separately excludes those agreeing
with an integer outside a finite set, including
finite-support exponents. Neither result determines
all intrinsic metabelian fibres or all congruence-cover
lifts.

\paragraph{Exact two-cover calculations behind the unit case.}
The following independent algebraic calculations preserve a useful
diagnostic for the rank-one cubic. They are not needed for the
general proof above. Put
\[
 E:y^2+y=x^3-1,\quad P=(1,0),\quad
 K=\Q(\delta),\quad\delta^2=-3,\quad\zeta=(\delta-1)/2.
\]
One has
$2P=(7,-19)$ and
$3P=(73/36,487/216)$, so $[3]P$ does not reduce to
$O$ at any rational prime at least five. A unit
exponent therefore gives unit values of every
boundary function outside the places above $2,3$.
For $f=\zeta-y$, the mod-two character lies in
\[
 K(\{2,3\},2)=
 \{1,-1,2,-2,\delta,-\delta,2\delta,-2\delta\}.
 \tag{A65.15}
\]
This follows from the Euclidean Eisenstein integer
ring, its unit group $\mu_6$, the inert prime $2$,
and the ramified prime generated by $\delta$.
For any representative $c$, the double cover
$w^2=(\zeta-y)/c$ has genus-two compactification
\[
 C_c:\quad V^2=U^6+D,\qquad
 U=x/w,\quad V=2c^2w-U^3,\quad D=4\delta c^3.
 \tag{A65.16}
\]
Its inverse has
$w=(V+U^3)/(2c^2)$, $x=Uw$, $y=\zeta-cw^2$.
The two degree-two maps
\[
 (X,Y)=(U^2,V),\qquad
 (X,Y)=(D/U^2,DV/U^3)
 \tag{A65.17}
\]
have targets $Y^2=X^3+D$ and $Y^2=X^3+D^2$.
Their differential pullbacks are $U\,dU/V$ and
$-dU/V$, so the Jacobian is isogenous to their
product. Since $D^2=-48c^6$, the second factor
is the fixed cubic $Y^2=X^3-48$, isomorphic
to $E$. After scaling, its map is the original
projection translated by $(0,\zeta)$.

These covers do not all have rank-zero Jacobians.
For $c=1$ the first factor $Y^2=X^3+4\delta$
has the point $(\zeta^2,2+\delta)$; its conjugate
gives a point on the $c=-1$ factor. At the places
$(q,\delta\bmod q)=(7,5),(13,6)$ the first curve
has respectively $4$ and $19$ points in its good
reductions. Prime-to-residue-characteristic torsion
injectivity at both places forces its entire
$K$-rational torsion group to be zero. The displayed
point is therefore nontorsion; the CM action
defined over $K$ gives rank at least two.
Thus these specific covers cannot be dismissed
by a rank-zero shortcut. The actual points
$P$ and $-P$ have classes $c=1$ and $c=-1$,
respectively.

There is also an exact rational norm comparison:
\[
 (\zeta-y)(\zeta^2-y)=x^3.
\]
Corestriction of the mod-two character gives the
$x$-character with square class $b=N_{K/\Q}(c)$,
which is $1$ for $c=\pm1,\pm2$ and $3$ for the
other four classes. Its rational double cover is
\[
 \mathcal C_b:\ V^2=4b^3w^6-3,\qquad
 x=bw^2,\quad y=(V-1)/2.
 \tag{A65.18}
\]
The two independent elliptic maps have coordinates
$(4b^3w^2,4b^3V)$ and $(-3/w^2,-3V/w^3)$,
and targets $Y^2=X^3-48b^6$ and
$Y^2=X^3+36b^3$. For $b=3$, the second target is
$E_{972}$, whose certified rank is zero in A.41;
the nonzero central value also gives finite
Tate--Shafarevich group by the CM theorem used in A.37.
Its good reduction orders at $5,7,13$ are $6,4,9$,
which force rational torsion to vanish; hence its
only rational point is $O$. The inverse fibre of
$O$ on $\mathcal C_3$ is $w=0,V^2=-3$, with no
rational point. Functorial finite abelian descent
and Stoll's finite-support theorem consequently
give $\mathcal C_3(\mathbb A_\Q)_\bullet^{\rm f\text{-}ab}
=\varnothing$.

For precision, a metabelian section lifts through
the twist selected by its global quadratic
character among these geometric-abelian covers, and then
supplies compact finite abelian descent: if $H$
is the geometric index-two subgroup and $\Delta$
the original geometric group, then
$\Delta'\subset H$ and $\Delta''\subset H'$.
The fixed-coset argument lifts the local geometry.
This gives a necessary cover condition, not a
converse lifting criterion. In the case $b=1$
both elliptic factors have positive rational rank
($E$ and $E_{36}$ have rank one by the certificates
of A.41 and A.43). The direct rational-unit proof
of A.65 settles the unit exponents without
assuming finite-descent excellence of this
positive-rank genus-two curve.

All identities, inverse maps, independent
differentials, exact finite-field counts and
cross-references to the saved rank certificates
in this diagnostic were executed in
\texttt{a\_main\_unit\_covers.py}; the complete
output is preserved. No numerical rank estimate
is used in place of those certificates.
